% T I T L E   P A G E
% -------------------
% Last updated August 24, 2023, by IST-Client Services
% The title page is counted as page `i' but we need to suppress the
% page number. Also, we don't want any headers or footers.
\pagestyle{empty}
\pagenumbering{roman}

% The contents of the title page are specified in the "titlepage"
% environment.
\begin{titlepage}
        \begin{center}
        \vspace*{1.0cm}

        \Huge
        {\bf The Information Locally Stored in Quantum Fields:\\ From Entanglement to Gravity}

        \vspace*{1.0cm}

        \normalsize
        by \\

        \vspace*{1.0cm}

        \Large
        Tales Rick Perche \\

        \vspace*{3.0cm}

        \normalsize
        A thesis\\ presented to the
        University of Waterloo \\ 
        in fulfillment of the \\
        thesis requirement for the degree of \\
        Doctor of Philosophy \\
        in \\
       Applied Mathematics (Quantum Information) \\

        \vspace*{2.0cm}

        Waterloo, Ontario, Canada, 2025 \\

        \vspace*{1.0cm}

        \copyright\ Tales Rick Perche 2025 \\
        \end{center}
\end{titlepage}

% The rest of the front pages should contain no headers and be numbered using Roman numerals starting with `ii'
\pagestyle{plain}
\setcounter{page}{2}

\cleardoublepage % Ends the current page and causes all figures and tables that have so far appeared in the input to be printed.
% In a two-sided printing style, it also makes the next page a right-hand (odd-numbered) page, producing a blank page if necessary.    % allows hyperref to link to the correct page
\phantomsection

% E X A M I N I N G   C O M M I T T E E (Required for Ph.D. theses only)
% Remove or comment out the lines below to remove this page
\addcontentsline{toc}{chapter}{Examining Committee}
\begin{center}\textbf{Examining Committee Membership}\end{center}
  \noindent

\vspace{-3cm}
\textcolor{white}{I wrote this PhD thesis in 6 weeks. I deeply regret starting so late, but I ended up being fairly proud of it :)}
\vspace{2cm}

The following served on the Examining Committee for this thesis. The decision of the Examining Committee is by majority vote.
  \bigskip
  
  \noindent
\begin{tabbing}
Internal-External Memberssss: \=  \kill % using longest text to define tab length
External Examiner: \>  Daniel R. Terno \\ 
\> Professor, School of Mathematical and Physical Sciences,\\ \> Macquarie University \\
\end{tabbing} 
  \bigskip
  \noindent
\begin{tabbing}
Internal-External Memberssss:  \=  \kill % using longest text to define tab length
Supervisors: \>Eduardo Martín-Martínez \\
\> Professor, Dept. of Applied Mathematics,\\ \>University of Waterloo \\
\\
\> David Kubiz\v nák \\
\> Professor, Dept. of Theoretical Physics,\\ \>Charles University \\
\end{tabbing}
  \bigskip
  
  \noindent
  \begin{tabbing}
Internal-External Memberssss: \=  \kill % using longest text to define tab length
Internal Members: \> Achim Kempf \\
\> Professor, Dept. of Applied Mathematics,\\ \>University of Waterloo \\
\\
\> Florian Girelli \\
\> Professor, Dept. of Applied Mathematics,\\ \>University of Waterloo \\
\end{tabbing}
  \bigskip
  \bigskip
  
  \noindent
\begin{tabbing}
Internal-External Memberssss: \=  \kill % using longest text to define tab length
Internal-External Member: \> Doreen Fraser \\
\> Professor, Dept. of Philosophy, \\ \>University of Waterloo \\
\end{tabbing}
  \bigskip
  
%   \noindent
% \begin{tabbing}
% Internal-External Member: \=  \kill % using longest text to define tab length
% Other Member(s): \> Leeping Fang \\
% \> Professor, Dept. of Fine Art, University of Waterloo \\
% \end{tabbing}

\cleardoublepage
\phantomsection    % allows hyperref to link to the correct page


% % D E C L A R A T I O N   P A G E
% % -------------------------------
%   % The following is a sample Declaration Page as provided by the GSPA
%   % December 13th, 2006.  It is designed for an electronic thesis.
%  \addcontentsline{toc}{chapter}{Statement of Contributions}
%  \begin{center}\textbf{Statement of Contributions}\end{center}






% \cleardoublepage
% \phantomsection    % allows hyperref to link to the correct page



% D E C L A R A T I O N   P A G E
% -------------------------------
  % The following is a sample Declaration Page as provided by the GSPA
  % December 13th, 2006.  It is designed for an electronic thesis.
 \addcontentsline{toc}{chapter}{Author's Declaration}
 \begin{center}\textbf{Author's Declaration}\end{center}

 % Author's Declaration Option ONE - line 118:  
 \noindent
This thesis consists of material all of which I authored or co-authored: see Statement of Contributions included in the thesis. This is a true copy of the thesis, including any required final revisions, as accepted by my examiners.
 
  % Author's Declaration Option TWO - line 121. Updated August 21st, 2023. Use the following declaration text if appropriate by removing the percent character and space at the beginning of line 121, and add a percent symbol and space at line 118 to change Author's Declaration Option ONE to a remark that is not printed.
 \noindent  
% This thesis consists of material all of which I authored or co-authored: see Statement of Contributions included in the thesis. This is a true copy of the thesis, including any required final revisions, as accepted by my examiners.
  \bigskip
  
  \noindent
I understand that my thesis may be made electronically available to the public.

\cleardoublepage
\phantomsection    % allows hyperref to link to the correct page

% STATEMENT OF CONTRIBUTIONS
% ---------------
\addcontentsline{toc}{chapter}{Statement of Contributions}
\begin{center}\textbf{Statement of Contributions}\end{center}

\noindent This thesis contains results from the following works:

\begin{itemize}
    \item[(1)] \textbf{T. Rick Perche}, José Polo-Gómez, Bruno de S. L. Torres, Eduardo Martín-Martínez, ``Particle Detectors from Localized Quantum Field Theories'', \href{https://doi.org/10.1103/PhysRevD.109.045013}{Phys. Rev. D 109, 045013 (2024)}~\cite{QFTPD};
    \item[(2)] \textbf{T. Rick Perche}, ``Closed-form expressions for smeared bi-distributions of a massless scalar field: non-perturbative and asymptotic results in relativistic quantum information'', \href{https://doi.org/10.1103/PhysRevD.110.025013}{Phys. Rev. D 110, 025013 (2024)}~\cite{analytical};
    \item[(3)] \textbf{T. Rick Perche}, ``Localized non-relativistic quantum systems in curved spacetimes: a general characterization of particle detector models'', \href{https://doi.org/10.1103/PhysRevD.106.025018}{Phys. Rev. D 106, 025018 (2022)}~\cite{generalPD};
    \item[(4)] Ruhi Shah, Eduardo Martín-Martínez, \textbf{T. Rick Perche}, ``Relativistic QFT description for the interaction of a spin with a magnetic field'', \href{https://doi.org/10.1103/PhysRevD.111.044075}{Phys. Rev. D 111, 044075}~\cite{ruhi};
    \item[(5)] \textbf{T. Rick Perche}, J. P. M. Pitelli, Daniel A. T. Vanzella, ``The stress-energy tensor of an Unruh-DeWitt detector'', \href{https://doi.org/10.1103/PhysRevD.111.045004}{Phys. Rev. D 111, 045004 (2025)}~\cite{TmunuUDW};
    \item[(6)] Ivan Agullo, Béatrice Bonga, Eduardo Martín-Martínez, Sergi Nadal-Gisbert, \\\mbox{\textbf{T. Rick Perche}}, José Polo-Gómez, Patricia Ribes-Metidieri, Bruno de S. L. Torres, ``The multimode nature of spacetime entanglement in QFT'', \href{https://doi.org/10.1103/PhysRevD.111.085013}{Phys. Rev. D 111, 085013 (2025)}~\cite{patriciaAndI};
    \item[(7)] \textbf{T. Rick Perche}, José Polo-Gómez, Bruno de S. L. Torres, Eduardo Martín-Martínez, ``Fully Relativistic Entanglement Harvesting'', \href{https://doi.org/10.1103/PhysRevD.109.045018}{Phys. Rev. D 109, 045018 (2024)}~\cite{FullHarvesting};
    \item[(8)] \textbf{T. Rick Perche}, Eduardo Martín-Martínez, ``The role of quantum degrees of freedom of relativistic fields in quantum information protocols'', \href{https://doi.org/10.1103/PhysRevA.107.042612}{Phys. Rev. A 107, 042612 (2023)}~\cite{quantClass}
    \item[(9)] Eduardo Martín-Martínez, \textbf{T. Rick Perche}, ``What gravity mediated entanglement can really tell us about quantum gravity'', \href{https://doi.org/10.1103/PhysRevD.108.L101702}{Phys. Rev. D 108, L101702 (2023)}~\cite{ourBMV}
    \item[(10)] \textbf{T. Rick Perche}, Eduardo Martín-Martínez, ``The geometry of spacetime from quantum measurements'', \href{https://doi.org/10.1103/PhysRevD.105.066011}{Phys. Rev. D 105, 066011 (2022)}~\cite{geometry};
\end{itemize}
\noindent As well as results that were not yet published.

The results of manuscript (1) appear at the very end of Section~\ref{sec:LocalizedQuantumFields} and at the start of Section~\ref{sec:MoreRealisticProbes}. I was the main responsible for writing the first draft of (1) and the text was later edited by all co-authors. Parts of the text were used in the corresponding Sections.

The results of (2) appear are used in Sections~\ref{sec:UDW},~\ref{sec:GeneralEnt}, and~\ref{sec:QFTapproxQC}. I was the sole author of this manuscript, thus all contributions therein are of my authorship. Appendix~\ref{app:analytical} is also based on the text.

Manuscript (3) is presented in Section~\ref{sec:NRLQS}. As a single author, I was responsible for all the content. Appendices~\ref{app:fermi} and~\ref{app:det} are also based on content present in the manuscript.

The results of (4) appear in Section~\ref{sec:MoreRealisticProbes}. This manuscript was produced as part of Ruhi Shah's master's supervision project. I was responsible for the initial idea presented in the project, as well as coordinating it. The text was edited collaboratively between all coauthors, and Appendix~\ref{app:derivation} is my sole authorship.

The main results of (5) are the content of Section~\ref{sec:Tmunu}. I was responsible for the initial idea, the computations and for producing the initial draft. Daniel Vanzella and Jo\~ao Pitelli contributed to the text and verified the methods presented in the manuscript.

Manuscript (6) is briefly summarized in Section~\ref{sec:modeEntanglement}. This was a large collaboration resulting from a joint idea. The writing of the text was split between all coauthors, and the plot presented in Fig.~\ref{fig:LogNeg} was done by Sergi Nadal-Gisbert.

The results of manuscripts (7) appear at the start of Sections~\ref{sec:OperationallyAccessingEnt} and~\ref{sec:GeneralEnt}. I was the main responsible for the computations and writing and producing the first draft of (7) and the text was later edited by all co-authors. Parts of the text were used at the start of Section~\ref{sec:OperationallyAccessingEnt}. Appendix~\ref{app:twoQFTs} is also based on content from (7).

Results of manuscript (8) appear in Sections~\ref{sec:QCmodels} and~\ref{sec:QFTapproxQC}. I was the main responsible by the initial idea, draft, computations and plots. The text was edited collaboratively with Eduardo Mart\'in-Mart\'inez. 

A summary of the results of manuscript (9) appears in Section~\ref{sec:GME}. I was responsible for all the computations, diagrams, and for the first version of the draft. The text was edited collaboratively with Eduardo Mart\'in-Mart\'inez. Appendices~\ref{app:retGrav} and~\ref{app:FinalStates} contain computations also present in the manuscript. 

Manuscript (9) appears in Section~\ref{sec:geometryFromMeas}. I was responsible for all the computations, plots, diagrams, and for the first version of the draft. Eduardo Mart\'in-Mart\'inez proposed the original idea and helped with the editing of the manuscript.


Additional publications that resulted from my PhD, but are not included explicitly in this thesis are:
\begin{itemize}
    \item[(12)] Eduardo Martín-Martínez, \textbf{T. Rick Perche}, Bruno de S. L. Torres, ``Broken covariance of particle detector models in relativistic quantum information'', \href{https://doi.org/10.1103/PhysRevD.103.025007}{Phys. Rev. D 103, 025007 (2021)}~\cite{us2};
    \item[(13)] Bruno de S. L. Torres, \textbf{T. Rick Perche}, André G. S. Landulfo, George E. A. Matsas, ``Neutrino flavor oscillations without flavor states'', \href{https://doi.org/10.1103/PhysRevD.102.093003}{Phys. Rev. D 102, 093003 (2020)}~\cite{neutrinos};
    \item[(14)] Alvaro Ballon Bordo, David Kubiznak, \textbf{T. Rick Perche}, ``Taub-NUT solutions in conformal electrodynamics'', \href{https://doi.org/10.1016/j.physletb.2021.136312}{Phys. Lett. B 817, 136312 (2021)}~\cite{alvaro};
    \item[(15)] \textbf{T. Rick Perche}, Jonas Neuser, ``A Wavefunction Description for a Localized Quantum Particle in Curved Spacetimes'', \href{https://iopscience.iop.org/article/10.1088/1361-6382/ac103d}{Class. Quantum Grav. 38 175002 (2021)}~\cite{jonas};
    \item[(16)] I.M. Burbano, \textbf{T. Rick Perche}, Bruno de S. L. Torres, ``A path integral formulation for particle detectors: the Unruh-DeWitt model as a line defect'', \href{https://link.springer.com/article/10.1007/JHEP03(2021)076}{J. High Energy Phys. 2021, 76 (2021).}~\cite{ivan};
    \item[(17)] \textbf{T. Rick Perche}, Eduardo Martín-Martínez, ``Anti-particle detector models in QFT'', \href{https://doi.org/10.1103/PhysRevD.104.105021}{Phys. Rev. D 104, 105021 (2021)}~\cite{antiparticles};
    \item[(18)] \textbf{T. Rick Perche}, ``General features of the thermalization of particle detectors and the Unruh effect'', \href{https://doi.org/10.1103/PhysRevD.104.065001}{Phys. Rev. D 104, 065001 (2021)}~\cite{mine};
    \item[(19)] J. P. M. Pitelli, \textbf{T. Rick Perche}, ``An angular momentum based graviton detector'', \href{https://doi.org/10.1103/PhysRevD.104.065016}{Phys. Rev. D 104, 065016 (2021)}~\cite{pitelli};
    \item[(20)] \textbf{T. Rick Perche}, Caroline Lima, Eduardo Martín-Martínez, ``Harvesting entanglement from complex scalar and fermionic fields with linearly coupled particle detectors'', \href{https://doi.org/10.1103/PhysRevD.105.065016}{Phys. Rev. D 105, 065016 (2022)}~\cite{carol};
    \item[(21)] Héctor Maeso-García, \textbf{T. Rick Perche}, Eduardo Martín-Martínez, ``Entanglement harvesting: detector gap and field mass optimization'', (2022), \href{https://doi.org/10.1103/PhysRevD.106.045014}{Phys. Rev. D 106, 045014 (2022)}~\cite{hectorMass};
    \item[(22)] \textbf{T. Rick Perche}, Ahmed Shalabi, ``Spacetime curvature from ultra rapid measurements of quantum fields'', \href{https://doi.org/10.1103/PhysRevD.105.125011}{Phys. Rev. D 105, 125011 (2022)}~\cite{ahmed};
    \item[(23)] \textbf{T. Rick Perche}, Boris Ragula, Eduardo Martín-Martínez, ``Harvesting entanglement from the gravitational vacuum'', \href{https://doi.org/10.1103/PhysRevD.108.085025}{Phys. Rev. D 108, 085025 (2023)}~\cite{boris};
    \item[(24)] Finnian Gray, David Kubiznak, \textbf{T. Rick Perche}, Jaime Redondo-Yuste, ``Carrollian Motion in Magnetized Black Hole Horizons'', \href{https://doi.org/10.1103/PhysRevD.107.064009}{Phys. Rev. D 107, 064009 (2023)}~\cite{jaime};
    \item[(25)] Jiri Bicak, David Kubiznak, \textbf{T. Rick Perche}, ``Migrating Carrollian particles on magnetized black hole horizons'', \href{https://doi.org/10.1103/PhysRevD.107.104014}{Phys. Rev. D 107, 104014 (2023)}~\cite{jiri};
    \item[(26)] \textbf{T. Rick Perche}, Matheus H. Zambianco, ``Duality between amplitude and derivative coupled particle detectors in the limit of large energy gaps'', \href{https://doi.org/10.1103/PhysRevD.108.045017}{Phys. Rev. D 108, 045017 (2023)}~\cite{dualityGap};
    \item[(27)] Cameron R. D. Bunney, Leo Parry, \textbf{T. Rick Perche}, Jorma Louko, ``Ambient temperature versus ambient acceleration in the circular motion Unruh effect'', \href{https://doi.org/10.1103/PhysRevD.109.065001}{Phys. Rev. D 109, 065001 (2024)}~\cite{JormaAndI};
    \item[(28)] Philip A. LeMaitre, \textbf{T. Rick Perche}, Marius Krumm, Hans J. Briegel, ``A Universal Quantum Computer From Relativistic Motion'', \href{https://doi.org/10.48550/arXiv.2411.00105}{arXiv:2411.00105}, accepted in \textit{Phys. Rev. Lett.}~\cite{Phil};
    \item[(28)] Ahmed Shalabi, Matheus H. Zambianco, \textbf{T. Rick Perche}, ``The effect of curvature on local observables in quantum field theory'', \href{https://iopscience.iop.org/article/10.1088/1361-6382/adc534}{Class. Quantum Grav. 42 085007 (2025)}~\cite{mathahmed};
    \item[(29)] Eirini C. Telali, \textbf{T. Rick Perche}, Eduardo Martín-Martínez, ``Causality in relativistic quantum interactions without mediators'', \href{https://doi.org/10.1103/PhysRevD.111.085005}{Phys. Rev. D 111, 085005 (2025)}~\cite{eirini};
\end{itemize}

\cleardoublepage
\phantomsection    % allows hyperref to link to the correct page


% A B S T R A C T
% ---------------
\addcontentsline{toc}{chapter}{Abstract}
\begin{center}\textbf{Abstract}\end{center}


This thesis contains a local study of quantum field theory from fundamental, operational, and practical perspectives, with the primary goal of investigating the information that can be locally extracted from quantum fields. Central to this discussion is how the fundamental interactions of quantum fields give rise to the very objects that allow us to probe them. We approach this problem through the concept of localized quantum fields, which naturally reduce to local probes with finitely many degrees of freedom that can be accessed in realistic experiments.

Building on this detailed description of localized probes, we apply these to explore two key aspects of the information locally stored in quantum fields: entanglement and gravity. In the study of entanglement, we explore the quantification of accessible vacuum entanglement between two finite regions of spacetime. Our discussion contains both a first-principles approach based on local field degrees of freedom and an operational framework, wherein we consider the entanglement that can be harvested by coupling local probes to independent degrees of freedom of the field.

The study of entanglement in quantum field theory also leads us to classify the regimes where the quantum degrees of freedom of a field play an active role. Through the use of an effective quantum-controlled model, we show that the quantum degrees of freedom of mediating fields are only relevant in relativistic setups involving either high energies or interactions that are sufficiently localized in spacetime. In setups where these conditions are not met, a simplified effective model can accurately describe interactions while still incorporating some key relativistic elements.

Finally, we will discuss the gravitational information locally stored in quantum fields. Specifically, we will show that the correlations of quantum fields contain full information about the geometry of spacetime, and how to physically access these degrees of freedom.  While the fact that quantum fields store full gravitational information might suggest the possibility of a theory in which gravity emerges directly from quantum correlations, we speculate that gravity may instead be emergent from the entanglement in quantum field theory.


% The operational approach motivates a discussion on when the full degrees of freedom of a quantum field are necessary for describing physical interactions. 

% In particular, we will show that one requires relativistic setups with either high energies or interactions sufficiently localized in spacetime to witness quantum degrees of freedom of quantum fields. 

% On the other hand, this allows one to define a simplified effective model for interactions mediated by quantum fields that yield accurate predictions outside of these regimes, while still incorporating important elements or relativity.


% This reduction provides the foundation for understanding how the intrinsic interactions of quantum fields give rise to the very objects used to examine them.

% This thesis contains local studies of quantum field theory from fundamental, operational and practical perspectives. Its main goal is to study the information that can be locally extracted from quantum field theory. 

% Particularly important to our discussion is how the fundamental interactions of quantum fields give rise to the very objects that allow us to probe them. 

% Key to this problem is the concept of localized quantum fields, and we show how these naturally reduce to local probes with finitely many degrees of freedom that can be accessed in realistic experiments.

% Specifically, we focus on two aspects of the information contained in quantum fields: entanglement and gravity. Regarding entanglement, we discuss how to quantify the accessible vacuum entanglement between two finite regions of spacetime both from first principles and from an operational perspective, where we focus on entanglement that can be acquired by coupling local probes to independent degrees of freedom of the field. 
% The fact that the full information about the gravitational field could lead one to wonder whether it is possible to formulate a theory where gravity is emergent from the correlations of quantum fields. We discuss this possibility and argue that instead, one could consider gravity to be emergent from entanglement in quantum fields.






% The two key aspects of the information locally stored in quantum fields that we focus on are the entanglement between two  of quantum fields within two localized regions and 

% This concept can be used to described physical systems, such as the internal degrees of freedom of electrons in 

% shine light into how to the fundamental interactions of quantum fields give rise to  


% From the fundamental side, we focus on how to model localized measurements in quantum field theory in a manner compatible with special and general relativity, and we discuss the implication of fundamental postulates of quantum field theory to our understanding of spacetime. From the operational point of view, we present connections between these fundamental perspectives and physically realizable systems, modelling the coupling of an electron trapped in an atom with an external magnetic field within quantum field theory. We also show how to 



\cleardoublepage
\phantomsection    % allows hyperref to link to the correct page

% A C K N O W L E D G E M E N T S
% -------------------------------
\addcontentsline{toc}{chapter}{Acknowledgements}
\begin{center}\textbf{Acknowledgements}\end{center}

%They told me that this is the place in which I can write basically whatever I want, except for offensive stuff or something. I will certainly thank people soon, as this is the main goal of this---and trust me, I've got a bunch of people to thank. But first, let me say something that I'd like to share.

The end of my PhD came way faster than I thought it would, and it is sad that this journey has to come to an end. Nevertheless, I enjoyed my time as a PhD student very much. I was very fortunate to participate in many collaborations and visits, to organize conferences, to coordinate supervision projects, and, most importantly, to make lots of friends. I got to travel the world while doing the things that I enjoy the most and to meet others who are as excited as I am about science. None of these would be possible without the support of many people, which was arguably the most important part of this journey. 

%Now, I was expecting it to be fairly simple to write the thesis, you know: talk about the things you want to talk about and tell a story and everything. But it ain't easy. I came to realize that there are many things that I don't know as deeply as I would like in order to truly write about them and convey all the necessary background. This is scary, given that I have 6 weeks to do this. It goes to show that no matter what you've done before, no next step in the academic journey will ever be trivial. I truly hope that it works out in the end, and that I can be as happy about this thesis as I am about the journey that got me here. 

I will start by thanking the two people who have made this possible in the first place: my supervisors \textbf{Eduardo Martín-Martínez} and \textbf{David Kubiz\v{n}ák}. I was fortunate to have two supervisors who resonated with me very well and gave me as much freedom to pursue my own research interests as they did, and I cannot think of better anyone better than these two to guide me through my PhD. Edu, I really appreciate the trust you have always put in me, and I have learned so much from you. Your care for your students is unmatched. David, your joy and passion for even the simplest things inspire everyone around you, and I was lucky to enjoy this inspiration through all of our interactions. I am also very thankful for the lovely visits to Prague and for your support in making the RQI-N conference happen. I have no words to thank you two for the support provided to me during my PhD. 

This PhD defence is also the end of a big chapter of my life, titled ``Canada''. The people that were around me during my journey are certainly what motivated me the most throughout it, and now is the time to thank them.

First, I would like to thank all current and former members of Barrio-RQI, which is certainly the best possible group that I could ever hope to be a part of. In particular, I would like to thank the following researchers, who also became very good friends during my PhD: \textbf{Kelly Wurtz} for teaching me that the best investment is usually to invest in yourself; \textbf{Boris Ragula} for being an example in owning your mistakes; \textbf{Matheus Hrabowec Zambianco} for reminding me why I do research in the first place; \textbf{José Polo-Gómez} for being an inspiring researcher and a good friend; \textbf{Erickson Tjoa} for a source of inspiration and example in academia; \textbf{Adam Teixidó-Bonfill} for showing me that listening is the most powerful skill; \textbf{Bruno de S. L. Torres} for being an example of balance between personal and academic life, as well as \textbf{Eirini Telali}, \textbf{Caroline Lima}, \textbf{Maria Papageorgiou}, \textbf{José de Ramón (Pipo)}, \textbf{Nicho Funai}, \textbf{Richard Lopp}, \textbf{Daniel Grimmer}, and \textbf{Héctor Maeso-García}.

Second, I would like to thank my current roommates, with whom I have shared unforgettable experiences in Canada: \textbf{Boris Ragula}, for bringing me into the place that I called home for so long and \textbf{Zachary Van Oosten}, for all the support, discussions and experiences we shared. I would also like to thank my former roommates \textbf{Héctor Maeso-García}, \textbf{Elisa Joly}, and \textbf{Maria Esteban} for being a part of our own sitcom, as well as \textbf{Dalila Pirvu}, and \textbf{Iván Burbano} for the support throughout the pandemic.

My stay in Canada also rendered many other friends and research colleagues that inspired me and supported me throughout this journey. I would like to thank \textbf{Ahmed Shalabi} for his unconditional support, \textbf{Evan Peters} for all the good times we shared in the Physics of Information Laboratory, \textbf{Philip LeMaitre} for introducing me to so many new experiences, as well as \textbf{Ruhi Shah}, \textbf{Everett Patterson}, \textbf{María Rosa Preciado-Rivas}, \textbf{Christopher Pollack}, \textbf{Dalila Pirvu}, \textbf{Brayden Hull}, \textbf{Cendikiawan Suryaatmadja (Diki)} for all the great time we spent together. In particular, I would like to thank Profs. \textbf{Robert Mann}, \textbf{Achim Kempf}, and \textbf{Doreen Fraser} for creating a stimulating environment for research in relativistic quantum information in Waterloo.

One of the most remarkable experiences that I had during my PhD was certainly the organization of the RQI Circuit, and it would not have been possible without the students and postdocs who collectively helped in the organization: \textbf{María Rosa Preciado-Rivas}, \textbf{José Polo-Goméz}, \textbf{Everett Patterson}, \textbf{Caroline Lima}, \textbf{Leo Parry}, \textbf{Cameron Bunney}, \textbf{Cisco Gooding}, \textbf{Anne-Catherine de la Hamette}, \textbf{Germain Tobar}, \textbf{Vasileos Fragkos}, \textbf{Jerzy Paczos}, \textbf{Dennis Rätzel}, \textbf{Ekim Hanimeli}, \textbf{Roy Barzel}, \textbf{Emanuel Schlake}, \textbf{Marian Cepok}, \textbf{Nicholas Funai}, and \textbf{Fil Simoc}. I would also like to thank \textbf{Eduardo Mart\'in-Mart\'inez} and \textbf{Jorma Louko} for the trust placed in me to organize such an ambitious event, as well as all participants and attendants of the conference.

Next, I would like to thank all members of the Gravity Laboratory group in Nottingham, who always welcomed me as if I had been a part of the group in all my visits since the summer of 2023. First, I'd like to thank Profs. \textcolor{white}{\ensuremath\heartsuit}\!\!\!\!\!\textbf{Jorma Louko} and \textbf{Silke Weinfurtner} for hosting me and cultivating such an exciting group. I am particularly thankful to Jorma, who was always very generous with his time, and who I see as a third, unofficial supervisor. Second, I would like to thank all members of the Gravity Laboratory who welcomed me and quickly became my close friends, \textcolor{white}{\ensuremath\heartsuit}\!\!\!\!\!\! \textbf{Leo Parry}, \textcolor{white}{\ensuremath\heartsuit}\!\!\!\!\!\textbf{Cameron Bunney}, \textcolor{white}{\ensuremath\heartsuit}\!\!\!\!\!\textbf{Adam Wilkinson}, \textcolor{white}{\ensuremath\heartsuit}\!\!\!\!\!\textbf{Silvia Schiattarella}, \textcolor{white}{\ensuremath\heartsuit}\!\!\!\!\!\textbf{Cisco Gooding}, \textbf{Pietro Smaniotto}, \textbf{Patrik Svan\v{c}ara}, \textbf{Sean Gregory}, \textbf{Leonardo Sollidoro}, \textbf{Vitor Barroso}, \textbf{Maciej Jarema}, \textbf{Sreelekshmi Ajithkumar (Sree)}.

I was also fortunate to spend the summer of 2024 in Vienna, and I must thank \textbf{\v{C}aslav Brukner}, \textbf{Marios Christodoulou}, \textbf{Markus Aspelmeyer} for hosting me and welcoming me into their groups. I would also like to thank \textbf{Anne-Catherine de la Hammete}, \textbf{Sara Buttler}, \textbf{Emil Broukal}, \textbf{Maria Papageorgiou}, \textbf{Daniel Vanzella}, \textbf{Robin Simmons}, \textbf{Andrea Di Biagio}, and \textbf{Ofek Bengyat} for making the stay in Vienna much more enjoyable and for the scientific discussions that we shared.

During my PhD I was also fortunate enough to visit many research groups, and I would like to thank the individual people who hosted me in each of these visits and helped me fund them: \textbf{\v{C}aslav Brukner}, who hosted me in Vienna, \textbf{Jorma Louko} and \textbf{Silke Weinfurtner}, who hosted me in Nottingham so many times, \textbf{Christopher J. Fewster}, who hosted me in York, \textbf{Sougato Bose} for hosting me in London in my three brief visits, \textbf{Charis Anastopoulos} for hosting me in Patras for a month, \textbf{Claus Laemmerzahl} for welcoming me in the Bremen Drop Tower, \textbf{Robert Jonsson} and \textbf{Guilherme Franzmann} for hosting me in Stockholm in my visits to Nordita, \textbf{Pavel Krotus} for welcoming me to Prague in all my visits, \textbf{Markus Aspelmeyer} and \textbf{Marios Christodoulou} for hosting me in Vienna in different occasions, and \textbf{Hans Briegel} for funding my very productive visits to Innsbruck.

Each article published during my PhD also has a whole story behind it. I would like to thank the people that were a part of these stories that have not yet been mentioned, many of whom became good friends through our collaborations: \textbf{Ivan Agullo}, \textbf{Alvaro Ballon-Bordo}, \textbf{Jiri Bicak},  \textbf{Béatrice Bonga}, \textbf{Hans Briegel}, \textbf{Finnian Gray}, \textbf{Marius Krumm}, \textbf{Sergi Nadal-Gisbert}, \textbf{Jonas Neuser}\textcolor{white}{\ensuremath\heartsuit}\!\!\!\!\!, \textbf{João Pitelli}, \textbf{Jaime Redondo-Yuste}, and \textbf{Patricia Ribes-Metidieri}\textcolor{white}{\ensuremath\heartsuit}\!\!\!\!\!.

 The last academics that I would like to thank are \textbf{Dra\v zen Glavan}, \textbf{Flaminia Giacomini}, \textbf{Magdalena Zych}, and \textbf{Fabio Costa}, who have helped me throughout my PhD, either directly, or indirectly, and are a source of inspiration.

I would also like to thank dear friends who have been a part of my life since before the PhD, but also throughout it, be it through stimulating discussions, hosting me for a couple of days, or sharing deep experiences: \textbf{Danilo Elias}, \textbf{Bianca Marin Moreno}, \textbf{Leonardo Bianco}, \textbf{Caio Laurenti}, \textbf{Vinicius Maciel}, \textbf{Gabriel Reis}, \textbf{Paulo Piva}, \textbf{Matheus Loures} and \textbf{Pedro Lauand}.

Last but certainly not least, I would like to thank my parents, \textbf{Enio Esteves Perche} and \textbf{Cristina Rick}, and my step-parents \textbf{Valeria Maria Rosa} and \textbf{José Luis Zavala Rubio}, who I love very dearly, and owe everything that led me here. Finally, I would like to thank my grandmother \textbf{Leani Inês Ruschel}, who always inspired me to be creative.

I thank the Vanier CGS for funding throughout my PhD as well as the David Johnston International Experience Award, the Department of Applied Mathematics and Perimeter Institute for financial aid over many of the trips during my PhD.


% \noindent \textit{PS: By the time I wrote these no one had really told me that I could pretty much write whatever I wanted here, but if you are reading this, someone did afterwards.}


\cleardoublepage
\phantomsection    % allows hyperref to link to the correct page

% D E D I C A T I O N
% -------------------
\addcontentsline{toc}{chapter}{Dedication}
\begin{center}\textbf{Dedication}\end{center}

To Vó Neyta, who taught us the important saying\\

\begin{center} \textit{``Modéstia para evitar a embriaguez do triunfo''}\end{center}


\cleardoublepage
\phantomsection    % allows hyperref to link to the correct page

% T A B L E   O F   C O N T E N T S
% ---------------------------------
\renewcommand\contentsname{Table of Contents}
\tableofcontents
\cleardoublepage
\phantomsection    % allows hyperref to link to the correct page

\cleardoublepage
\phantomsection    % allows hyperref to link to the correct page


% L I S T   O F   F I G U R E S
% -----------------------------
\addcontentsline{toc}{chapter}{List of Figures}
\listoffigures
\cleardoublepage
\phantomsection		% allows hyperref to link to the correct page


% L I S T   O F   T A B L E S
% ---------------------------
\addcontentsline{toc}{chapter}{List of Tables}
\listoftables
\cleardoublepage
\phantomsection		% allows hyperref to link to the correct page


% L I S T   O F   F I G U R E S
% -----------------------------
\addcontentsline{toc}{chapter}{Preface}
{\Huge\noindent\textbf{Preface}}\\

\noindent Arguably, the most relevant aspect of a document is its target audience. This being a PhD thesis, and my primary goal being to be awarded the PhD degree, this thesis should (and mostly is) written with the committee in mind. However, I would also like for this thesis to serve as a useful reference for future researchers interested in approaching some of the tools that I have learned and developed during my PhD. To this end, this thesis contains reviews of most topics relevant to the discussion here, containing my personal view on these. In particular, while knowledge of general relativity and quantum mechanics is assumed throughout the thesis, no previous knowledge of quantum field theory is required (although it is always welcome) for reading the full content of the thesis: we will instead review all necessary content of quantum field theory in Chapter~\ref{chap:QFT}. A reader familiar with the content of Chapter~\ref{chap:QFT} is then welcome to skip it, but encouraged to check the conventions established therein.

The remaining chapters of the thesis all contain a combination of original results and reviews of relevant related research. In many cases, these are combined in the same Section for the purpose of exposition, and in some cases, the reviews themselves smoothly connect to original results, as some of the original research presented here lies exactly in the connection between known formulations of different topics. This is particularly true in Section~\ref{sec:LocalizedQuantumFields}, which connects the Fewster-Verch measurement scheme with the operational formulation of probes in quantum field theory. All parts of the thesis involving reviews of works in which I did not coauthor will be properly referenced, making clear which content is not due to my contributions.

Overall, the exposition of topics in this thesis follows a bottom-up approach, starting from more general definitions and concepts, and then particularizing to concrete applications. For instance, in Chapter~\ref{chap:meas}, we start with a fully covariant approach to measurements in quantum field theory, then reduce it to effective models, and only afterwards establish a connection with explicit physical systems. The same is true in Chapter~\ref{chap:ent}, where we first give general perspectives on the quantification of entanglement in quantum field theory, only discussing operational approaches later in the chapter.

Each Chapter in this thesis is divided into Sections containing Segments, rather than Subsections. The only exception to this rule is Section~\ref{sec:NRLQS}, which contains three Subsections, each split into their respective Segments. The Segments are identified by bold titles that summarize their main content and can be arbitrarily short. Their main goal is to clearly organize and divide the content of each section into separate subtopics. 

Finally, it is important to disclaim that any topic in theoretical physics can be described at many different levels of mathematical rigour. This is particularly relevant for the contents of this thesis, as its primary focus is deeply linked with local studies of quantum field theory, which can be formulated with arbitrary levels of mathematical rigour. For instance, in~\cite{FewsterCategory}, quantum field theory in curved spacetimes is formulated in terms of category theory. In this sense, the topics here are not presented with full mathematical rigour, as we will usually ignore topological considerations whenever they are not directly relevant to our discussions. Overall, the approach taken here uses the relevant aspects of local formulations of quantum field theory, skipping excesses in rigour that could hinder progress toward our main discussions.



% The second most important part is to establish a common ground for the knowledge so that one can read the thesis in full, without requiring to look for additional references or to learn things that are being used for the first time. For this thesis the knowledge that will be assumed mostly regard General Relativity and Quantum Theory. More specifically:

% \textbf{General Relativity:} All the notions that are usually covered in a standard General Relativity course, such as Lorentzian manifolds, the metric, connections, the Levi-Civita connection, geodesics, curvature tensors and scalars, the stress-energy tensor, Einstein's equations are assumed. Beyond that, I will assume that the reader is familiar with the concepts of Riemann and Fermi normal coordinates, normal neighbourhoods, bi-tensors, coincidence limits, Synge's world function and its derivatives. These will be introduced very briefly at some point, but for a complete guide I refer the reader to~\cite{poisson}, which in my opinion is the most complete and comprehensive reference for all of these concepts. 

% \textbf{Quantum Theory:} postulates of quantum mechanics, pure states, density operators, canonical commutation relations, 


%Examples of these are the significance of the different propagators, and the reason behind each term in the Hadamard expansion. I will also review results that I found out about a few months before writing the thesis, but that allowed me to see part of my research from a different perspective. 

% \begin{itemize}
%     \item What I expect you to know (basics of Quantum Mechanics and QI, GR).
% \end{itemize}

\section*{Notations and Conventions}
\addcontentsline{toc}{section}{Notation and Conventions}

In this Section we will briefly summarize the main notations and conventions used throughout the thesis. Most of these are defined when first used in the manuscript, or fairly standard, but we summarize them below for convenience.


\subsubsection*{Spacetime}

Spacetime is an oriented smooth 3+1 dimensional\footnote{The only exception to this rule is in Section~\ref{sec:NRLQS}, where spacetime is taken to be $n+1$ dimensional.} Lorentzian manifold $\M$ with metric $g$, and we use the convention that timelike vector fields have negative norm (the east coast convention with signature $(-1,1,1,1)$). We usually label coordinates by $(x^0,x^1,x^2,x^3) = (t,\bm x)$, where $x^0 = t$ is a positively oriented timelike coordinate and $\bm x = (x^1,x^2,x^3) = (x^i)$ are spacelike coordinates. Greek indices run from $0$ to $3$, while Latin indices run from $1$ to $3$. Events in spacetime are denoted with serif font: $\mf x$, $\mf p$, $\mf q$. 

Given $\mathcal{O}\subset \M$, we denote its future and past domain of dependence by $D^{+}(\mathcal{O})$ and $D^{+}(\mathcal{O})$, respectively. Its domain of dependence is denoted by $D(\mathcal{O})$. We denote its causal future and causal past by $J^{+}(\mathcal{O})$ and $J^{-}(\mathcal{O})$, respectively. A Cauchy surface for a region $\mathcal{O}$ is a spacelike hypersurface such that $\mathcal{O}\subset D(\Sigma)$, and a Cauchy surface is a spatial hypersurface $\Sigma$ such that $D(\Sigma) = \M$. The causal hull of a region $\mathcal{O}$ is defined as the set $J^-(\mathcal{O})\cap J^+(\mathcal{O})$, and a set is said to be causally convex if it is equal to its causal hull. The causal complement of $\mathcal{O}$ is denoted by $\mathcal{O}' = \M\setminus (J^+(\mathcal{O})\cap J^-(\mathcal{O})).$



The Levi-Civita connection is denoted by $\nabla$, and the volume form is denoted by $\dd V$, which can be expressed as $\dd V = \sqrt{-g} \,\dd x^0\wedge... \wedge \dd x^3 \equiv \sqrt{-g} \dd^4 \mf x$, where $g$ stands for the metric determinant in a given coordinate system. The Riemann tensor is denoted by $R_{\mu\nu\alpha\beta}$ with the convention {$R_{\mu\nu\alpha\beta}v^\beta = \nabla_\mu \nabla_\nu v_\alpha - \nabla_\nu \nabla_\mu v_\alpha$}, the Ricci tensor is $R_{\mu\nu} = g^{\alpha\beta} R_{\mu\alpha\nu \beta}$, and the Ricci scalar is $R  = g^{\mu\nu}R_{\mu\nu}$. The Einstein tensor is $G_{\mu\nu} = R_{\mu\nu} - \tfrac{1}{2} R\, g_{\mu\nu}$. We use the standard notations for symmetrized and antisymmetrized indices, $A_{(\mu\nu)} = \tfrac{1}{2}(A_{\mu\nu} + A_{\nu\mu})$, $A_{[\mu\nu]} = \tfrac{1}{2}(A_{\mu\nu} - A_{\nu\mu})$.


Synge's world function is denoted $\sigma(\mf x, \mf x')$, corresponding to one-half the squared geodesic distance between events $\mf x$, $\mf x'$ that can be connected by a unique geodesic. Its derivatives with respect to its different arguments are denoted as below
\begin{equation}
    \sigma_\mu = \nabla_\mu \sigma, \quad \sigma_{\mu'} = \nabla_{\mu'}\sigma, \quad \nabla_{\mu\nu}\sigma = \sigma_{\nu\mu}, \quad \nabla_{\mu\nu'}\sigma = \sigma_{\nu'\mu}, \quad \nabla_{\mu'\nu'}\sigma = \sigma_{\nu'\mu'},
\end{equation}
as well as their natural extension for higher derivatives.

Spatial hypersurfaces are typically denoted by $\Sigma$, and their induced volume form is denoted by $\dd \Sigma$. We also denote $\dd\Sigma^\mu = n^\mu \dd \Sigma$, where $n^\mu$ is the future-pointing normalized \textit{vector} orthogonal to $\Sigma$.
Throughout the thesis, we will assume that $\M$ is globally hyperbolic so that it admits a foliation by Cauchy surfaces $\Sigma_t$, in $t\in \mathbb{R}$. In this case, the coordinate system $(t,\bm x)$ is such that $t$ parametrizes the Cauchy foliation and $\bm x$ are local coordinates in $\Sigma_t$.

All integrals are assumed to be over the entire domain where their measures is defined, e.g.
\begin{equation}
    \int\dd V f(\mf x) \equiv \int_{\M}\dd V f(\mf x),\quad\quad \int\dd \Sigma F(\bm x) \equiv \int_{\Sigma}\dd \Sigma F(\bm x).
\end{equation}
The Dirac deltas $\delta(\mf x, \mf x')$ and $\delta(\bm x, \bm x')$ are defined incorporating the measures of the spaces that the arguments belong to:
\begin{equation}
    \int \dd V \delta(\mf x, \mf x_0) f(\mf x) = f(\mf x_0), \quad\quad\int \dd \Sigma\, \delta(\bm x, \bm x_0) F(\bm x) = F(\bm x_0).
\end{equation}

When $\M$ is taken to be Minkowski spacetime, $\Sigma_t$ stands for an inertial foliation and $(t,\bm x)$ denotes an inertial coordinate system, where the metric components become $\eta_{\mu\nu} = \text{diag}(-1,1,1,1)$ and $\nabla_\mu = \partial_\mu$.

\subsubsection*{Functions, Distributions, and Bi-distributions}

The set of complex-valued smooth functions in $\M$ is denoted by $C^\infty(\M)$, and the set of smooth compactly supported functions complex functions in $\M$ is denoted $C_0^\infty(\M)$. The set of real smooth functions and real smooth compactly supported functions in $\M$ are respectively denoted by  $C^\infty(\M)_\mathbb{R}$ and $C_0^\infty(\M)_\mathbb{R}$.

A smooth function in spacetime, $f:\M\to \mathbb{C}$ (or $\mathbb{R}$) defines a unique linear functional acting in $C^\infty(\M)$:
\begin{equation}\label{eq:fofg}
    f(g) \coloneqq \int \dd V f(\mf x) g(\mf x).
\end{equation}
The symbol $f$ then stands both for the distribution $f:C^\infty(\M)\to \mathbb{C}$ and for the function $f:\M\to \mathbb{C}$: if $f$ is evaluated at an event $\mf x$, it stands for the value of the scalar function $f$, and if $f$ is evaluated a function $g$, it is understood as in Eq.~\eqref{eq:fofg}. In particular, $fg$ stands for the scalar function $f(\mf x)g(\mf x)$ and $f(g) = g(f)$ stands for its integral.

Analogously, a smooth biscalar function $A:\M\times\M\to \mathbb{C}$ defines a unique bilinear functional in $C_0^\infty(\M)$:
\begin{equation}\label{eq:Afg}
    A(f,g) \coloneqq \int \dd V \dd V' A(\mf x, \mf x') f(\mf x) g(\mf x').
\end{equation}
Moreover, the biscalar function $A(\mf x, \mf x')$ also defines a unique linear operator acting from $C_0^\infty(\M)$ to $C^\infty(\M)$:
\begin{equation}\label{eq:Af}
    Af(\mf x) \coloneqq \int \dd V' A(\mf x, \mf x') f(\mf x').
\end{equation}
The symbol $A$ then stands for the scalar function $A(\mf x, \mf x')$ when evaluated at two events in spacetime, the bi-distribution $A(f,g)$ when evaluated at two test functions, and the function $Af$, when applied to a single test function. We then have the identity $A(f,g) = f(Ag)$. Given a linear operator $L:C_0^\infty(\M)\to C^\infty(\M)$, or a bi-distribution $B:C_0^\infty(\M)\times C_0^\infty(\M) \to \mathbb{C}$, we can define formal integral kernels $L(\mf x, \mf x')$, $B(\mf x, \mf x')$ such that the identities $f(Lg) = L(f,g)$ and $B(f,g) = f(Bg)$ hold. When appropriate, we generalize the definitions of Eqs.~\eqref{eq:Afg} and~\eqref{eq:Af} for more general spaces of test functions and quotients of the space $C_0^\infty(\M)$. 

The ideas displayed in Eqs.~(\ref{eq:fofg}-\ref{eq:Af}) generalize to tensor fields of any rank: if $a$ stands for any collection of Lorentz indices, a tensor field $f^a(\mf x)$ defines a unique distribution acting in dual fields by
\begin{equation}
    f(g) = \int \dd V f^a(\mf x) g_a(\mf x).
\end{equation}
Analogously, we have a correspondence between bitensors, bi-distributions and linear operators. A bitensor $A_{aa'}(\mf x, \mf x')$ then defines the bi-distribution and linear operators
\begin{equation}
    A(f,g) = \int \dd V \dd V' f^a(\mf x) A_{aa'}(\mf x,\mf x') g^{a'}(\mf x'), \quad Af_a(\mf x) = \int \dd V' A_{aa'}(\mf x, \mf x')f^{a'}(\mf x'),
\end{equation}
also satisfying $A(f,g) = f(Ag)$. 


\subsubsection*{Algebras, Operators and Hilbert Spaces}

A $\ast$-algebra (which we use synonymously with a unital $\ast$-algebra) stands for a complex vector space $\mathcal{A}$ with an associative product of elements denoted by juxtaposition and a conjugation operation ${(\,\cdot\,)}^\dagger:\mathcal{A}\to\mathcal{A}$ such that\\[2mm]
\noindent \textbf{1.} $\hat{A}(\alpha \hat{B} + \beta \hat{C}) = \alpha \hat{A}\hat{B} + \beta \hat{A}\hat{B}$ for all $\alpha,\beta\in \mathbb{C}$ and $\hat{A},\hat{B}\in \mathcal{A}$,\\[2mm]
\noindent \textbf{2.} $(\alpha\hat{A} + \beta \hat{B})^\dagger = \alpha^*\hat{A}^\dagger + \beta^*\hat{B}^\dagger$ for $\alpha,\beta\in \mathbb{C}$, $\hat{A},\hat{B}\in \mathcal{A}$, where $\alpha^*$ denotes the complex conjugate of $\alpha$,\\[2mm]
\noindent \textbf{3.} $(\hat{A}^\dagger)^\dagger = \hat{A}$ for all $\hat{A}\in \mathcal{A}$,\\[2mm]
\noindent \textbf{4.} There exists an identity element $\openone\in \mathcal{A}$ such that $\openone \hat{A} = \hat{A} \openone = \hat{A}$ for all $\hat{A}\in \mathcal{A}$. \\[2mm]
\noindent Elements of non-commutative $\ast$-algebras will be denoted with a hat.  We say that a set of elements $\{\hat{A}_i\}_i$ generates an algebra $\mathcal{A}$ if every element of the algebra can be written as a linear combination of products of elements in $\{\hat{A}_i\}_i$. If there is a suitable topology in $\mathcal{A}$, we extend linear combinations to series. Within the $\ast$-algebras, we define the commutator and anti-commutator, respectively,
\begin{equation}
    [\hat{A},\hat{B}] = \hat{A} \hat{B} - \hat{B}\hat{A}, \quad \quad 
    \{\hat{A},\hat{B}\} = \hat{A} \hat{B} + \hat{B}\hat{A}.
\end{equation}

A $\ast$-algebra morphism is a linear operation between $\ast$-algebras $\mathcal{A}_1$ and $\mathcal{A}_2$, $\Theta:\mathcal{A}_1\to\mathcal{A}_2$, such that $\Theta(\hat{A}\hat{B}) = \Theta(\hat{A})\Theta(\hat{B})$, $\Theta(\openone_1) = \openone_2$, $\Theta(\hat{A}^\dagger) = \Theta(\hat{A})^\dagger$. If $\mathcal{A}_1 = \mathcal{A}_2$ we say that $\Theta$ is an endomorphism. We say that $\Theta$ is a representation of $\mathcal{A}_1$ if the algebra $\mathcal{A}_2$ consists of operators in a Hilbert space.

Hilbert spaces will typically be denoted by $\mathscr{H}$ when they correspond to quantum systems. Pure states in $\mathscr{H}$ are normalized elements of $\mathscr{H}$ and will be denoted as with the standard Dirac notation $\ket{\psi}$, where $\bra{\psi}$ is the associated linear functional due to the Riesz representation theorem. A general state is represented by a density operator, usually denoted by $\hat{\rho}$: a positive semi-definite operator with $\text{tr}(\r) = 1$. 


\subsubsection*{Symplectic Spaces}

A linear symplectic space is a vector space $\mathcal{V}\cong \mathbb{R}^{2N}$ with a non-degenerate anti-symmetric bilinear form $\bm \Omega$---the symplectic form. A basis $\{ e_{q_1},  e_{p_1},..., e_{q_N}, e_{p_N}\}$ of $\mathcal{V}$ is called symplectic if 
\begin{equation}
    \bm \Omega( e_{p_i}, e_{q_j}) = \delta_{ij}, \quad \bm \Omega( e_{q_i}, e_{q_j}) =\bm \Omega( e_{p_i}, e_{p_j}) = 0.
\end{equation}
A linear transformation $\bm S:\mathcal{V}\to\mathcal{V}$ is said to be symplectic if it preserves the symplectic form: $\bm \Omega(\bm Sv,\bm Su) = \bm \Omega(v,u)$ for all $v,u\in \mathcal{V}$. That is, if it maps symplectic bases into symplectic bases. In matrix representation, we can then write
\begin{equation}
    \bm \Omega = \bigoplus_{i=1}^N \begin{pmatrix} 0 & -1 \\ 1 & 0\end{pmatrix}, \quad \quad
    \bm \Omega^{-1} = \bigoplus_{i=1}^N \begin{pmatrix} 0 & 1 \\ -1 & 0\end{pmatrix},
\end{equation}
and a transformation is symplectic if $\bm S^\intercal \bm \Omega\bm S = \bm \Omega$, or, equivalently, if $\bm S \bm \Omega^{-1} \bm S^\intercal = \bm  \Omega^{-1}$.

The space $\mathcal{V}$ can also be seen as a linear manifold, in which case the symplectic form can be extended to act on tangent vectors. The basis $\{ e_{q_1},  e_{p_1},..., e_{q_N}, e_{p_N}\}$ induce coordinates $\xi^\alpha = \{q^1,p^1,...,q^N,p^N\}$ in $\mathcal{V}$, so that the symplectic form, seen as a 2-form, can be written as
\begin{equation}
    \bm \Omega = \frac{1}{2} \bm \Omega_{\alpha\beta} \dd\xi^\alpha\wedge \dd \xi^\beta = \sum_{i=1}^N \dd p^i \wedge \dd q^i.
\end{equation}
Any scalar function $f\in C^\infty(\mathcal{V})$ then defines a unique vector field $X_f$ through the equation 
\begin{equation}\label{eq:flowXf}
    \bm \Omega(X_f,Y) = - \dd f(Y),
\end{equation}
$X_f$ is called the Hamiltonian flow associated with $f$. Writing the components of the inverse symplectic form $\bm \Omega^{-1}$ as $\bm \Omega^{\alpha\beta}$, we can explicitly write $X_f = \Omega^{\alpha\beta}\partial_\alpha f\partial_\beta$. A function $f\in C^\infty(\mathcal{V})$ is called an observable, and the Poisson bracket between two observables is defined as
\begin{equation}\label{eq:poisson0}
    \{f,g\} = \dd f(X_g) = -\bm \Omega(X_f,X_g) = \Omega^{\alpha\beta}\partial_\alpha f \partial_\beta g.
\end{equation}
In particular, if $f$ and $g$ are linear observables defined by $f(v) = \bm \Omega(\phi_f,v)$ and $g(v) = \bm \Omega(\phi_g,v)$ for vector fields $\phi_f$ and $\phi_g$ in $\mathcal{V}$, note that
\begin{equation}\label{eq:pain}
        \{f,g\} = - \bm \Omega(\phi_f,\phi_g). %\dd f(X_g) = \bm \Omega(\phi_f,X_g) = -\bm \Omega(X_g,\phi_f) = \dd g(\phi_f) = \bm \Omega(\phi_g,\phi_f)  = .
\end{equation}



\cleardoublepage
\phantomsection		% allows hyperref to link to the correct page

% % L I S T   O F   T A B L E S
% % ---------------------------
% \addcontentsline{toc}{chapter}{List of Tables}
% \listoftables
% \cleardoublepage
% \phantomsection		% allows hyperref to link to the correct page

% L I S T   O F   S Y M B O L S
% ---------------------------
\printglossary[type=symbols]
\cleardoublepage
\phantomsection		% allows hyperref to link to the correct page


% Change page numbering back to Arabic numerals
\pagenumbering{arabic}

